% Propietario conservado: /Users/ruben/Documents/ChatGPT/jueces y controles/REVISION_ARGUMENTAL_APERTURAS_CIERRES_20260915/II/manuscript_es/sections/02b_arquitectura_tpk.tex
% Fragmento acumulativo para insertar después de la presentación de la
% actualización TPK en nucleo.tex. No sustituye nucleo, extension,
% generacion ni registro_k. Procedencia detallada en el MD compañero.
\subsection{Arquitectura operatoria del transporte, la memoria y la prolongación}
\label{ii:ii:tpk-arquitectura}

El núcleo topológico de fase organiza la evolución de las dos evaluaciones
de APP sobre historias orientadas. Su estructura reúne un soporte de caminos,
reglas locales de actualización, relaciones de prolongación, registros
temporales y aplicaciones de lectura. Cada componente tiene un dominio
propio. Esta jerarquía permite conservar, en una misma construcción, la
posición observable, el cociente de transporte y la memoria de las
fronteras atravesadas. La composición de la ecuación~\eqref{ii:eq:actualizacion}
adquiere así una realización precisa en cada uno de esos niveles.

\subsubsection{Caminos bivariados y transporte levantado}

Sea \(\mathsf P_{\rm APP}\) la categoría libre del grafo cardinal sobre
\(X_9=(\mathbb Z/9\mathbb Z)^2\). Sus objetos son las posiciones de APP;
sus morfismos son las palabras cardinales, con composición por
concatenación. La categoría
\[
 \mathsf P_{\rm APP}^{(2)}
 =\mathsf P_{\rm APP}\times\mathsf P_{\rm APP}
\]
conserva ordenadamente el camino de la hoja aditiva y el de la hoja
multiplicativa. Ambos caminos recorren el mismo soporte aritmético y
mantienen diferenciadas sus evaluaciones.

En la carta \(I_9=\{0,\ldots,8\}\), los generadores son
\[
 \nu_N=(-1,0),\quad\nu_E=(0,1),\quad
 \nu_S=(1,0),\quad\nu_O=(0,-1),
 \qquad T_d(x)=[x+\nu_d]_9.
\]
Para \(w=d_1\cdots d_r\), el levantamiento al recubridor entero conserva
\[
 \widetilde x_k=\widetilde x_0+\sum_{j=1}^{k}\nu_{d_j},
 \qquad x_k=[\widetilde x_k]_9.
\]
Cuando \(x_r=x_0\), la clase de enrollamiento es
\begin{equation}
 \operatorname{wind}(w)
 =\frac{1}{9}(\#S-\#N,\#E-\#O)\in\mathbb Z^2.
 \label{ii:ii:tpk-enrollamiento}
\end{equation}
La posición final y el enrollamiento son dos lecturas diferentes de la
misma ruta. La segunda registra el desplazamiento del levantamiento aun
cuando la primera haya retornado.

La inversión de caminos tiene dominio y codominio determinados:
\[
 \operatorname{rev}:\operatorname{Hom}(x,y)
 \longrightarrow\operatorname{Hom}(y,x),\qquad
 d_1\cdots d_r\longmapsto\overline d_r\cdots\overline d_1,
\]
donde \(\overline N=S\), \(\overline E=O\). Satisface
\(\operatorname{rev}^2=I\), invierte el orden de composición y cambia el
signo de~\eqref{ii:ii:tpk-enrollamiento}. Ésta es la operación precisa que
debe transportarse al comparar orientaciones conjugadas de una ruta.

\subsubsection{Estado completo y conservación de las actualizaciones}

La descripción de \(\widehat X\) se organiza sobre las posiciones y
direcciones de los dos cursores. Sobre cada configuración observable
\(q\), la fibra aritmética se escribe
\begin{equation}
 \mathcal F_q=
 \mathsf O_q\times\mathsf H_q\times\mathsf C_q\times
 \mathsf S_q\times\mathsf B_q\times\symup{\Lambda}_q.
 \label{ii:ii:tpk-fibra}
\end{equation}
Las coordenadas \(\mathsf O_q\) conservan orientación y hoja;
\(\mathsf H_q\), los residuos y cocientes de las evaluaciones;
\(\mathsf C_q\), los acarreos compatibles;
\(\mathsf S_q\), la firma de frontera;
\(\mathsf B_q\), el prefijo y sus cilindros; y
\(\symup{\Lambda}_q\), el registro ordenado de transporte. La fibra está
restringida por las identidades aritméticas, de incidencia y de
truncamiento correspondientes. El estado completo conserva además la ruta
que lo ha producido y el origen desde el que se leen sus fases.

La selección de modo tiene tipo
\[
 M_{\rm ph}:D_9\longrightarrow\{-1,0,+1\},\qquad
 \mathsf{Sel}_t:\widehat X\longrightarrow\widehat X^{\rm sel}_t.
\]
La segunda aplicación incorpora al estado el valor de la primera y
activa la hoja correspondiente. El transporte levantado
\[
 \mathsf{Tra}_t:\widehat X^{\rm sel}_t
 \longrightarrow\widehat X^{\rm tra}_t
\]
aplica \(T_d\) al cursor activo, concatena su letra de ruta y conserva la
hoja seleccionada. Finalmente,
\[
 \mathsf{Upd}_t:\widehat X^{\rm tra}_t\longrightarrow\widehat X
\]
actualiza las evaluaciones y los registros. En esta notación,
\(\mathsf{Upd}_t\) es la actualización \(\mathsf{Mem}_t\) de
\eqref{ii:eq:actualizacion}, explicitada en sus componentes.

Cada evaluación entera \(a\) se conserva mediante
\(a=\rho_9(a)+9q_9(a)\); cada evaluación ternaria, mediante su resto
orientado y su acarreo. Para una concatenación \(r_2r_1\), el transporte
de los acarreos satisface la ley cocíclica
\begin{equation}
 k_{\rm tr}(r_2r_1)
 =k_{\rm tr}(r_2)+\mathsf C(r_2)k_{\rm tr}(r_1),
 \label{ii:ii:tpk-cociclo-transporte}
\end{equation}
donde \(\mathsf C(r_2)\) es la acción de transporte de la segunda ruta
sobre la coordenada de acarreo de la primera. En el registro ordenado,
la actualización se realiza por adjunción del evento producido:
\[
 \Lambda_{t+1}=\Lambda_t\Vert
 (\text{hoja},d_t,\text{residuo},\text{cociente},
   \text{acarreo},\text{orientación},\text{frontera})_t.
\]
La conservación tiene aquí un significado dinámico: las contribuciones
anteriores permanecen recuperables dentro del registro que se prolonga.
Sus coordenadas pueden cambiar al transportar el estado.

\subsubsection{Acción afín sobre la memoria de transporte}

Para una trayectoria \(r:q\to q'\), las aplicaciones
\[
 \mathsf H(r):\mathsf H_q\to\mathsf H_{q'},\qquad
 \mathsf C(r):\mathsf C_q\to\mathsf C_{q'}
\]
respetan identidad y concatenación; \(\mathsf C_q\) es el grupo
abeliano de memoria de transporte y \(\mathsf C(r)\) su homomorfismo.
El incremento \(k_{\rm tr}(r)\in\mathsf C_{q'}\), con
\(k_{\rm tr}(\operatorname{id}_q)=0\), es el cociclo de
\eqref{ii:ii:tpk-cociclo-transporte}. Este símbolo distingue el incremento
de transporte de las coordenadas racionales del registro dodecafásico.

En una carta con datos \((h,c,\lambda)\), la actualización tiene la
forma explícita
\begin{equation}
 h'=\mathsf H(r)h,\qquad
 c'=\mathsf C(r)c+k_{\rm tr}(r),\qquad
 \lambda'=r\circ\lambda.
 \label{ii:ii:tpk-accion-afin}
\end{equation}
La firma se evalúa mediante
\(\operatorname{Sig}_{q'}(h',c',\lambda')\). Su residencia en la
fibra \(\mathsf S_{q'}\) conserva su dependencia de esos datos.

\begin{proposition}[Composición de la memoria afín]
\label{ii:ii:tpk-composicion-afin}
La actualización por \(r_1\), seguida de la actualización por
\(r_2\), coincide con la actualización por \(r_2r_1\).
\end{proposition}
\begin{proof}
Dos actualizaciones de la memoria dan
\[
 c''=\mathsf C(r_2)\mathsf C(r_1)c
       +\mathsf C(r_2)k_{\rm tr}(r_1)+k_{\rm tr}(r_2).
\]
La compatibilidad de \(\mathsf C\) con la concatenación y la
identidad~\eqref{ii:ii:tpk-cociclo-transporte} convierten esta expresión
en \(\mathsf C(r_2r_1)c+k_{\rm tr}(r_2r_1)\).
La coordenada \(h\) satisface la misma compatibilidad por
functorialidad de \(\mathsf H\); la trayectoria la satisface por
asociatividad. La firma coincide al evaluarse sobre los mismos datos.
\end{proof}

La carta espacial de cada cursor proporciona una realización elemental
de esa ley. Para \(x\in I_9^2\) y una dirección cardinal \(d\),
\[
 a_d(x)=\frac{x+\nu_d-T_d(x)}9\in\mathbb Z^2,
 \qquad x'=T_d(x),\qquad c'=c+a_d(x).
\]
Aquí \(\mathsf C(r)=I\), y \(k_{\rm tr}(r)\) suma los
cruces enteros de las fronteras. Para una ruta de longitud \(n\),
la suma telescópica da
\[
 x_n=x_0+\sum_{j=1}^n\nu_{d_j}
             -9\sum_{j=1}^n a_{d_j}(x_{j-1}).
\]
En un lazo, la última suma coincide con
\(\operatorname{wind}(r)\). La posición residual retorna mientras
el levantamiento conserva el número y orientación de sus cruces.
Las componentes restantes de la memoria mantienen sus aplicaciones
propias de transporte y sus relaciones de prolongación.

\subsubsection{Cronología de los dos cursores y retorno observable}

El calendario finito vive en \(\Theta_{108}=\mathbb Z/108\mathbb Z\).
Para \(0\leq\widetilde\theta<108\), ponemos
\[
 r(\theta)=1+(\widetilde\theta\bmod9),\qquad
 \beta(\theta)=\left[\left\lfloor\frac{\widetilde\theta}{9}
 \right\rfloor\right]_{12},\qquad
 \varepsilon_9(r)=M_{\rm ph}(r).
\]
La presentación~\eqref{ii:tpk-int:qobs} agrupa posición y dirección por
cursor. En esta ampliación se agrupan primero las dos posiciones y
después las dos direcciones. Ambas cartas se identifican mediante la
permutación de factores
\begin{equation}
 \begin{aligned}
 \mathfrak p_{\rm obs}:\ (X_9\times\mathcal D)^2\times\Theta_{108}
 &\longrightarrow X_9^2\times\mathcal D^2\times\Theta_{108},\\
 (x_+,d_+,x_\times,d_\times,\theta)
 &\longmapsto(x_+,x_\times,d_+,d_\times,\theta),
 \end{aligned}
 \label{ii:ii:permutacion-coordenadas-observables}
\end{equation}
con \(\mathcal D=\{N,E,S,O\}\). La inversa repone la dirección junto
a su cursor. Si los superíndices \({\rm agr}\) y \({\rm sep}\) distinguen
estas dos escrituras del sucesor observable, entonces
\[
 \mathfrak p_{\rm obs}F_{\rm obs}^{\rm agr}
 =F_{\rm obs}^{\rm sep}\mathfrak p_{\rm obs}.
\]
En efecto, el selector depende de la misma fase, el vector cardinal se
aplica al mismo cursor y la inversión de direcciones ocurre en el mismo
instante; únicamente cambia el orden de las coordenadas. Conservamos
\(\mathcal Q_{\rm obs}\) y \(F_{\rm obs}\) para el espacio y la
aplicación así identificados. Esta permutación de presentación deja
inalterados los eventos y las coordenadas de memoria de la historia.

El espacio observable es
\[
 \mathcal Q_{\rm obs}
 =X_9^2\times\{N,E,S,O\}^2\times\Theta_{108}.
\]
Escribamos un estado como
\(q=(x_+,x_\times,d_+,d_\times,\theta)\). Primero se evalúa la hoja
activa en la posición actual. Después se actualizan las posiciones por
\begin{equation}
 (x_+',x_\times')=
 \begin{cases}
 (T_{d_+}x_+,x_\times),&\varepsilon_9(r)=+1,\\
 (x_+,T_{d_\times}x_\times),&\varepsilon_9(r)=-1,\\
 (x_+,x_\times),&\varepsilon_9(r)=0.
 \end{cases}
 \label{ii:ii:tpk-fobs-posiciones}
\end{equation}
Si \(\widetilde\theta+1\) es múltiplo de \(27\), se sustituyen
\((d_+,d_\times)\) por \((\overline d_+,\overline d_\times)\).
En todos los casos se incrementa \(\theta\) módulo \(108\). Estas reglas
definen el sucesor
\(F_{\rm obs}:\mathcal Q_{\rm obs}\to\mathcal Q_{\rm obs}\).
La realización de emisión de la sección~\ref{ii:sec:extension} aplica
exactamente este orden de evaluación, transporte e inversión.

\begin{proposition}[Retorno observable y conservación de la historia]
\label{ii:ii:tpk-retorno-observable}
La aplicación \(F_{\rm obs}\) es una permutación. El operador
\(\mathcal K_{27}=F_{\rm obs}^{27}\) satisface
\(\mathcal K_{27}^{4}=I\). Una historia completa de ciento ocho
actualizaciones conserva además los eventos y el acarreo acumulados
durante esa vuelta.
\end{proposition}
\begin{proof}
Conocida la fase posterior se recupera la fase anterior. La regla indica
si se produjo una inversión de direcciones; deshacerla y aplicar el
transporte cardinal inverso al cursor activo recupera el estado previo.
En cada tramo de veintisiete transiciones hay nueve activaciones de cada
hoja. Desde un extremo de tramo, cada cursor realiza nueve desplazamientos
en la misma dirección, cuya suma es cero en \(X_9\); después invierte
su orientación. Dos tramos restauran las orientaciones y cuatro restauran
también el calendario. La identidad se transporta a cualquier fase por
la invertibilidad de \(F_{\rm obs}\). El registro completo conserva la
concatenación de los ciento ocho eventos y su contador, que se han
actualizado en cada transición.
\end{proof}

La periodicidad de esta permutación pertenece al cociente observable. En
la historia completa, el contador entero de ventanas aumenta en doce
durante el mismo intervalo. La sucesión exacta
\eqref{ii:eq:extension108} expresa la reducción del contador a su coordenada
dodecafásica y conserva su defecto de aditividad.


% BEGIN OWNER /Users/ruben/Documents/ChatGPT/jueces y controles/REVISION_ARGUMENTAL_APERTURAS_CIERRES_20260915/II/manuscript_es/sections/02c_transporte_electronico.tex
\subsubsection{Transporte electrónico bidireccional y registro aritmético}
\label{ii:ii:transporte-electronico}

La representación central utiliza una celda compartida por las dos
evaluaciones APP. Sus coordenadas pertenecen a \(D_9=\{1,\ldots,9\}\),
y la inversión \((r,c)\mapsto(10-r,10-c)\) tiene el punto fijo único
\((5,5)\). El estado visible se escribe \((r,c,h,v,n)\), con
\(h,v\in\{-1,1\}\) e índice cronológico entero \(n\geq0\).
El selector trítico repite \((+1,-1,0)\). En la hoja aditiva se lee
\(r+c\) y se traslada la columna en la dirección \(h\); en la
multiplicativa se lee \(rc\) y se traslada la fila en la dirección
\(v\). El umbral conserva la posición. Ambas orientaciones se
invierten al completar cada intervalo de veintisiete actualizaciones.

Para expresar la lectura y el acarreo con el mismo convenio positivo,
definimos
\[
 [z]^+_9=1+((z-1)\bmod9),\qquad
 q_9(z)=\frac{z-[z]^+_9}{9}.
\]
La evaluación precede al desplazamiento. Si \(a_n\) es el entero
leído, el evento conserva \((a_n,[a_n]^+_9,q_9(a_n))\), la hoja,
la posición anterior, la orientación y el cruce de frontera. El
símbolo trítico es \([a_n]^+_9\bmod3\). El levantamiento espacial
retiene, por ejemplo, el entero
\(q_9(c+h)\) cuando \(c\) se sustituye por \([c+h]^+_9\).

La representación de celda compartida y la representación de dos
cursores de \eqref{ii:ii:tpk-fobs-posiciones} tienen estados visibles
diferentes. Sus registros quedan relacionados con el estado TPK por
sus mapas de representación respectivos. Para la trayectoria central
se utiliza aquí, de forma íntegra, la actualización sobre la celda
compartida que produce la palabra electrónica.

\begin{proposition}[Actualización invertible con cocientes conservados]
\label{ii:ii:electron-inversa-cocientes}
Añadamos al estado visible dos registros enteros \(Q_+,Q_\times\)
y dos números de enrollamiento \(W_r,W_c\). Cada lectura activa
incrementa el registro de su hoja en \(q_9(a_n)\); cada desplazamiento
incrementa el enrollamiento correspondiente en su cruce de frontera.
Esta actualización posee una inversa explícita sobre su imagen.
\end{proposition}
\begin{proof}
El índice posterior determina el modo anterior y si acaba de ocurrir
una inversión de orientaciones. Se deshace primero esa inversión.
Si el modo era aditivo, se recupera la columna previa como
\([c-h]^+_9\); si era multiplicativo, se recupera la fila como
\([r-v]^+_9\). El modo neutro conserva ambas coordenadas. La celda
recuperada determina nuevamente el entero leído, su residuo, su
cociente y el cruce de frontera. Restando esos incrementos de
\(Q_+,Q_\times,W_r,W_c\) e invirtiendo el índice cronológico se
recupera exactamente el estado precedente. La composición de estas
inversas recupera una historia finita completa. La concatenación de
eventos conserva además el orden de las lecturas.
\end{proof}

Los registros \(Q_+,Q_\times\) son una publicación aditiva explícita
de la historia de cocientes. La construcción se integra en la fibra
de memoria junto con sus demás coordenadas. Su ley de composición,
para historias concatenables \(u,v\), es
\[
 Q_s(uv)=Q_s(u)+Q_s(v),\qquad s\in\{+,\times\},
\]
donde la segunda historia comienza en el estado final de la primera.
La inversión algebraica de una actualización resta sus incrementos;
recorrer posteriormente un trayecto espacial de vuelta registra
nuevas evaluaciones. Esa distinción conserva simultáneamente la
reversibilidad del transporte y la duración de la historia.

\begin{proposition}[Retorno central y acumulación exacta]
\label{ii:ii:electron-memoria-retorno}
Desde \((r,c,h,v,n)=(5,5,1,1,0)\) y registros iniciales nulos, la
proyección espacial y sus orientaciones retornan después de cincuenta y cuatro
actualizaciones. En ese intervalo,
\[
 (\Delta W_r,\Delta W_c)=(0,0),\qquad
 (\Delta Q_+,\Delta Q_\times)=(10,46).
\]
Después de \(54N\) actualizaciones, para todo entero \(N\geq0\),
se tiene \((Q_+,Q_\times)=(10N,46N)\), mientras las coordenadas
espaciales y sus orientaciones retornan. El índice entero toma el valor
\(54N\); el calendario dodecafásico módulo \(108\) retorna para
\(N\) par.
\end{proposition}
\begin{proof}
Cada par de modos activos desplaza una vez cada coordenada. En
veintisiete actualizaciones ambas recorren una vuelta de nueve
posiciones y retornan a la diagonal central; después se invierten
\(h,v\). El bloque siguiente recorre la vuelta opuesta. Los cruces
espaciales se cancelan entre ambos bloques.

En cada bloque las lecturas aditivas recorren \(2k\),
\(1\leq k\leq9\), cuyo cociente positivo suma cinco. Las lecturas
multiplicativas del bloque positivo recorren
\(k[k+1]^+_9\). Su suma entera y su suma de residuos son
\[
 \sum_{k=1}^9 k[k+1]^+_9=249,\qquad
 \sum_{k=1}^9[k[k+1]^+_9]^+_9=42.
\]
La segunda identidad resulta de la lista
\((2,6,3,2,3,6,2,9,9)\). Por división euclídea,
la suma de cocientes es \((249-42)/9=23\). La sustitución
\(k\mapsto[k-1]^+_9\) muestra que el bloque negativo contiene
el mismo multiconjunto de productos. Ambos bloques producen, por
tanto, diez y cuarenta y seis unidades de los registros respectivos.
La proyección espacial y sus orientaciones vuelven a sus valores
iniciales. Las reglas de activación y de inversión se repiten al
desplazar el índice en \(54\), y los incrementos no dependen de los
registros acumulados. La ley aditiva
anterior prueba por inducción la fórmula para todo \(N\).
\end{proof}

La lectura trítica y la orientación por ventanas proporcionan, sobre
la misma trayectoria,
\[
 \gamma_e=\bigl((100000220)^3(120200000)^3\bigr)^2,\qquad
 \tau_e=(+,+,+,-,-,-,+,+,+,-,-,-).
\]
Las potencias representan concatenaciones. Los lectores de
discrepancias y de ventanas desarrollados en
\eqref{ii:eq:el-kd} y \eqref{ii:eq:el-komega} calculan ambos
\((80,54,6)\). La evaluación se obtiene de la trayectoria antes
de incorporar \(\alpha\), la sección de acción o una masa.
La memoria de cocientes conserva información adicional a ese triple:
en ciento ocho actualizaciones registra \((20,92)\), con
enrollamiento espacial neto nulo. Cada publicación mantiene así el
tipo del observable que la produce.

% END OWNER



% BEGIN OWNER /Users/ruben/Documents/ChatGPT/jueces y controles/REVISION_ARGUMENTAL_APERTURAS_CIERRES_20260915/II/manuscript_es/sections/02d_bidireccion_memoria.tex
% Desarrollo separado, 9 de septiembre de 2026.
% Amplía el lector producto de centro_electronico_registros.tex.
% No identifica la media vuelta espacial con una conjugación física.
\subsubsection{Bidireccionalidad del lector producto y memoria de cocientes}
\label{ii:bim09:seccion}

La lectura de producto conserva cuatro coordenadas aritméticas antes de
publicar una firma. Para $x=(i,j)\in D_9^2$, $D_9=\{1,\ldots,9\}$,
se definen
\[
 N(x)=ij,\qquad r(x)=1+((ij-1)\bmod9),\qquad
 q(x)=\frac{N(x)-r(x)}9,
\]
y el lector discreto
\[
 \ell(1,2,4,8,7,5)=(0,1,2,3,4,5),\qquad
 \ell(3)=\ell(6)=\ell(9)=0.
\]
Así, $N=r+9q$ conserva simultáneamente producto entero, residuo positivo
y cociente aritmético. La función $f_\ell(x)=\ell(r(x))$ es la
coordenada utilizada por la firma regional. Las funciones $N,r,q$ y
$f_\ell$ tienen dominio posicional; la dirección de transporte permanece
en la semilla orientada.

\paragraph{Avance, media vuelta y lectura anterior al transporte.}

Sea $\mathscr X_\times=D_9^2\times\{N,E,S,O\}$. Los vectores cardinales
son $\nu_N=(-1,0)$, $\nu_E=(0,1)$, $\nu_S=(1,0)$,
$\nu_O=(0,-1)$. La suma posicional se reduce a representantes positivos
módulo nueve. Definimos
\[
 P(x,d)=(x+\nu_d,d),\qquad H(x,d)=(x,\bar d),
 \qquad \nu_{\bar d}=-\nu_d.
\]
Estos operadores satisfacen $P^9=I$, $H^2=I$ y $HPH=P^{-1}$.

En las primeras cincuenta y cuatro transiciones TPK, el canal producto
se activa en las fases $2,5,8$ de cada ventana nonádica. Cada activación
lee la posición anterior al avance. Tras las transiciones $27$ y $54$
se aplica $H$. Por tanto, las dieciocho posiciones leídas son, para una
semilla $s$,
\[
 P^0s,P^1s,\ldots,P^8s,\quad
 HP^9s,\,P H P^9s,\ldots,P^8H P^9s,
\]
entendiendo que la lectura posicional omite únicamente la dirección.
Para cualquier función posicional $f$, pongamos $f_n=f(P^ns)$, con
índices módulo nueve. La palabra de lecturas crudas es
\begin{equation}
 \mathcal R_f(s)=
 (f_0,f_1,\ldots,f_8,\ f_0,f_8,f_7,\ldots,f_1).
 \label{ii:bim09:palabra}
\end{equation}
Las seis sumas de ventana, cada una de tres activaciones, son
\begin{align}
 A_j^f(s)&=f_{3j}+f_{3j+1}+f_{3j+2},&&0\leq j<3,\label{ii:bim09:ida}\\
 A_{3+j}^f(s)&=f_{-3j}+f_{-3j-1}+f_{-3j-2},&&0\leq j<3.
 \label{ii:bim09:vuelta}
\end{align}
Para $f=f_\ell$, la firma es
$E_{\rm sig}(s)=([A_0^f]_{10},\ldots,[A_5^f]_{10})$.
El cociente decimal de cada acumulador es
$c_j^{10}=\lfloor A_j^f/10\rfloor$ y conserva
$A_j^f=[A_j^f]_{10}+10c_j^{10}$.

\begin{proposition}[Media vuelta y permutación de semiciclos]
\label{ii:bim09:media-vuelta}
Para toda semilla orientada y toda función posicional $f$,
\[
 \mathcal R_f(Hs)=\operatorname{rot}_9\mathcal R_f(s),\qquad
 A_j^f(Hs)=A_{j+3\bmod6}^f(s).
\]
La misma permutación transporta los residuos y los cocientes decimales
de los acumuladores. Se aplica, en particular, a $N,r,q$ y $f_\ell$.
\end{proposition}
\begin{proof}
La relación $P^nH=HP^{-n}$ y la igualdad $f(Hs)=f(s)$ dan
$f(P^nHs)=f(P^{-n}s)=f_{-n}$. Sustituir en
\eqref{ii:bim09:palabra} intercambia sus dos grupos de nueve entradas.
La suma en grupos de tres produce la identidad de ventanas. La división
euclídea por diez se efectúa componente a componente y respeta la
permutación.
\end{proof}

El avance también tiene una ley exacta sobre los acumuladores:
\begin{align}
 A_j^f(Ps)-A_j^f(s)&=f_{3j+3}-f_{3j},&&0\leq j<3,
 \label{ii:bim09:avance-ida}\\
 A_{3+j}^f(Ps)-A_{3+j}^f(s)&=f_{1-3j}-f_{-2-3j},&&0\leq j<3.
 \label{ii:bim09:avance-vuelta}
\end{align}
La prueba consiste en cancelar los dos sumandos comunes de cada ventana.
El transporte de una firma requiere, por tanto, las evaluaciones de
frontera que aparecen en estas diferencias, además de las sumas ya
publicadas.

\paragraph{Inversión cronológica y términos de frontera.}

Sea $\gamma=(x_0,\ldots,x_n)$ una trayectoria registrada y sea
$\gamma^\dagger=(x_n,\ldots,x_0)$ su inversión cronológica, con cada
arista orientada en sentido opuesto. Para una función posicional $f$,
la lectura anterior a cada avance es
\[
 C_f(\gamma)=\sum_{t=0}^{n-1}f(x_t).
\]
Entonces
\begin{equation}
 C_f(\gamma^\dagger)=C_f(\gamma)+f(x_n)-f(x_0).
 \label{ii:bim09:frontera-inversion}
\end{equation}
Si $n=3m$ y $A_j^f$ agrupa tres aristas consecutivas,
\begin{equation}
 A_j^f(\gamma^\dagger)=A_{m-1-j}^f(\gamma)
 +f(x_{3(m-j)})-f(x_{3(m-j-1)}).
 \label{ii:bim09:frontera-ventana}
\end{equation}
Ambas fórmulas se siguen de que la trayectoria inversa lee los extremos
de llegada de las aristas originales. Los sumandos interiores coinciden;
la diferencia conserva los extremos.

Para un observable de arista antisimétrico $a(x,y)=-a(y,x)$, su palabra
sí se transforma mediante
\[
 \mathcal R_a(\gamma^\dagger)
 =-\operatorname{rev}\mathcal R_a(\gamma).
\]
Esta identidad actúa sobre aristas orientadas. Las identidades
\eqref{ii:bim09:frontera-inversion}--\eqref{ii:bim09:frontera-ventana}
actúan sobre lecturas posicionales. La media vuelta $H$ cambia la
dirección inicial y conserva el orden del calendario; la inversión
cronológica invierte el registro completo. Sus dominios y sus reglas de
lectura quedan especificados por separado.

\paragraph{Cociclo de ocupación y memoria conservada en ida y vuelta.}

Para una arista cardinal $x\to y$, con representantes positivos
$\widetilde x,\widetilde y$, el cruce del recubrimiento es
\[
 b(x,y)=\frac{\widetilde x+\nu_d-\widetilde y}{9}\in\mathbb Z^2.
\]
Se cumple $b(y,x)=-b(x,y)$. En la categoría libre de caminos
registrados, definimos el lector aditivo
\begin{equation}
 \mathfrak m(\gamma)=
 \left(\sum_{t=0}^{n-1}b(x_t,x_{t+1}),\ C_q(\gamma),\ n\right).
 \label{ii:bim09:memoria-aditiva}
\end{equation}
La concatenación satisface
$\mathfrak m(\gamma_2\circ\gamma_1)=
\mathfrak m(\gamma_1)+\mathfrak m(\gamma_2)$.
Es un cociclo aditivo sobre trayectorias registradas: el registro conserva
las aristas de ida y de vuelta como eventos diferentes.

Por \eqref{ii:bim09:frontera-inversion},
\begin{equation}
 \mathfrak m(\gamma^\dagger\circ\gamma)=
 \bigl(0,\ 2C_q(\gamma)+q(x_n)-q(x_0),\ 2n\bigr),
 \label{ii:bim09:memoria-retorno}
\end{equation}
donde $\gamma^\dagger\circ\gamma$ denota la concatenación temporal
de la ida seguida de la vuelta. La suma orientada de cruces se cancela;
la ocupación aritmética y la longitud permanecen. Este lector no
desciende al cociente que identifica un recorrido seguido de su inversa
con la trayectoria vacía, puesto que ese cociente elimina los eventos
que la segunda y la tercera coordenadas conservan.

Hay también una lectura par puramente fronteriza:
\[
 V_b(\gamma)=\sum_{t=0}^{n-1}
 b(x_t,x_{t+1})b(x_t,x_{t+1})^{\mathsf T},\qquad
 V_b(\gamma^\dagger\circ\gamma)=2V_b(\gamma).
\]
La lectura impar y la lectura par registran propiedades distintas de
las mismas aristas; ambas están calculadas antes de olvidar el camino.

\paragraph{Aplicación exacta a la fibra de firma \texorpdfstring{$555555$}{555555}.}

La enumeración finita de $\mathscr X_\times$ contiene $324$ semillas,
$26$ firmas y $24$ semillas con firma $555555$. La partición estable
por $P,H$ refina esta fibra en tres clases de ocho elementos. Pueden
etiquetarse por su firma después de un avance:
\[
 \mathscr C_{177},\quad\mathscr C_{717},\quad\mathscr C_{771}.
\]
La acción de media vuelta es
\[
 H\mathscr C_{177}=\mathscr C_{177},\qquad
 H\mathscr C_{717}=\mathscr C_{771},\qquad
 H\mathscr C_{771}=\mathscr C_{717}.
\]
El certificado finito adjunto verifica estas identidades sobre todos los
representantes. En las veinticuatro semillas, los seis acumuladores
crudos de $f_\ell$ valen exactamente cinco. Por tanto, los cocientes
decimales $c_j^{10}$ se anulan en esta fibra. La información adicional
se observa en las posiciones, las evaluaciones enteras y los cocientes
aritméticos, que tienen lecturas diferentes de esa firma.

La coordenada transversal fija $a$ del recorrido cardinal vale $4$ o
$5$ en esas semillas. Durante una vuelta de nueve avances se recorre
una vez cada valor de la coordenada móvil $u\in D_9$. Puesto que
$\gcd(a,9)=1$, los residuos positivos $r(au)$ permutan $D_9$.
De aquí se obtiene, sin utilizar una constante física,
\[
 \sum_{u=1}^9au=45a,\qquad
 \sum_{u=1}^9r(au)=45,\qquad
 \sum_{u=1}^9q(au)=5(a-1).
\]
La vuelta posterior recorre las mismas posiciones en sentido opuesto.
En las dieciocho activaciones,
\begin{equation}
 \sum N=90a,\qquad\sum r=90,\qquad
 \sum q=10(a-1)=
 \begin{cases}30,&a=4,\\40,&a=5.\end{cases}
 \label{ii:bim09:ocupacion-30-40}
\end{equation}
El producto entero acumulado vale, respectivamente, $360$ o $450$.
El enrollamiento neto es cero; la suma de cocientes y las dieciocho
incidencias conservan la historia de la lectura. Cada una de las tres
clases de firma contiene representantes de ambos valores transversales.
La misma firma y la misma clase de transporte de ese cociente finito
pueden, por tanto, coexistir con ocupaciones aritméticas distintas.

Estos resultados conectan la memoria conservada con operaciones
concretas: permutación de semiciclos, diferencias de frontera,
cocientes aritméticos y registro de incidencias. La identificación
físico-matemática de cualquier lector posterior conserva sus mapas
propios; las identidades aquí demostradas pertenecen a la dinámica
interna del canal producto.

% END OWNER


% BEGIN OWNER /Users/ruben/Documents/ChatGPT/jueces y controles/REVISION_ARGUMENTAL_APERTURAS_CIERRES_20260915/II/manuscript_es/sections/02e_estado_producto.tex
% Fragmento de inserción breve. Desarrollo y procedencia en la misma carpeta.
\subsubsection{Memoria aritmética del canal multiplicativo}
\label{ii:prd-arit:seccion}

La conservación de la firma regional se prolonga a la evaluación
aritmética completa de APP. Para una posición orientada
\(x=(i,j,d)\in\{1,\ldots,9\}^2\times\{N,E,S,O\}\), sea
\(\chi(x)=(a,b,\epsilon)\), donde \(a\) es la coordenada fija,
\(b\) la coordenada móvil y \(\epsilon\) su orientación:
\[
 \chi(i,j,N)=(j,i,-1),\quad \chi(i,j,S)=(j,i,+1),\qquad
 \chi(i,j,E)=(i,j,+1),\quad \chi(i,j,O)=(i,j,-1).
\]
La etiqueta de eje \(\iota\in\{i,j\}\) conserva cuál de las
dos coordenadas originales se mueve. Escribiendo
\(\langle b\rangle_9=1+((b-1)\bmod9)\), los operadores y la
evaluación son
\begin{align}
 \overline P(a,b,\epsilon)&=(a,\langle b+\epsilon\rangle_9,\epsilon),
 &\overline H(a,b,\epsilon)&=(a,b,-\epsilon),\label{ii:prd-arit:PH}\\
 n(a,b)&=ab,&r(a,b)&=\langle ab\rangle_9,\qquad
 q(a,b)=\frac{ab-r(a,b)}9.\label{ii:prd-arit:rq}
\end{align}

\begin{proposition}[Representación exacta del transporte aritmético]
\label{ii:prd-arit:minimal}
Se cumplen \(\chi P=\overline P\chi\) y
\(\chi H=\overline H\chi\). La representación \(\chi\)
tiene \(162\) estados y es la menos fina de las representaciones
estables bajo \(P,H\) que conservan el producto entero o,
equivalentemente, el par \((r,q)\). Conservar además \(\iota\)
recupera las \(324\) posiciones orientadas originales.
\end{proposition}
\begin{proof}
El avance modifica sólo \(b\), y la media vuelta sólo
\(\epsilon\), lo que prueba las dos identidades de transporte.
Las nueve observaciones \(n_k=n(\overline P^k\chi(x))\),
\(0\leq k<9\), recorren \(a,2a,\ldots,9a\). Así se recuperan
\(a=\min_k n_k\), \(b=n_0/a\) y el signo mediante
\((n_1/a-n_0/a)\bmod9\in\{1,8\}\). Toda equivalencia estable
que conserve \(n\) conserva estas nueve observaciones y, por
tanto, distingue todos los estados \((a,b,\epsilon)\).
La identidad \(n=r+9q\) prueba la equivalencia de las dos
evaluaciones. Cada fibra de \(\chi\) contiene las dos rutas
transpuestas: \((i,j,d)\mapsto(j,i,d^\top)\), con
\(N^\top=O,O^\top=N,E^\top=S,S^\top=E\). La etiqueta de eje distingue
esas dos preimágenes y permite reconstruir posición y dirección.
\end{proof}

El selector TRIT activa el producto en \(t\equiv2,5,8\pmod9\),
antes de cada avance. La media vuelta actúa al finalizar cada
\(27\) instantes. Si \(y_{t-1}\) denota el estado aritmético
previo, el incremento y su memoria acumulada quedan determinados por
\begin{equation}
 \eta_q(t)=\mathbf1_{\{t\equiv2\ (3)\}}q(y_{t-1}),\qquad
 \lambda_q(t)=\lambda_q(t-1)+\eta_q(t),\quad \lambda_q(0)=0.
 \label{ii:prd-arit:incremento}
\end{equation}
Se cumple \(0\leq q(a,b)\leq a-1\). Para invertir una
actualización se deshace primero la media vuelta que corresponda,
después el avance, y finalmente se resta el cociente recuperado del
estado previo. El registro retiene fase, eje, orden de eventos y
memoria; \(\lambda_q\) es un cociclo aditivo de este transporte.

Cada bloque de \(27\) contiene nueve avances y recorre una vez
la coordenada móvil. Por ello, el retorno de posición y orientación
en \(54\), y de fase en \(108\), coexiste con
\begin{equation}
 \lambda_q(108N)=NQ_\times^+(a),\qquad
 Q_\times^+(a)=\frac{180a-4\sum_{b=1}^9r(a,b)}9,
 \qquad N\geq0.
 \label{ii:prd-arit:retorno}
\end{equation}
En efecto, cada bloque suma los productos \(a,2a,\ldots,9a\);
sumar \(n=r+9q\) en los cuatro bloques da la fórmula. Si
\(a\) es primo con \(9\), los residuos suman \(45\), y
\(Q_\times^+(a)=20(a-1)\). En la ventana \(j\), contando desde
\(j=0\), el signo relativo del calendario es
\(\sigma_j=(-1)^{\lfloor j/3\rfloor}\), y la orientación
efectiva es \(\epsilon_j=\epsilon_0\sigma_j\).
La carga firmada por \(\sigma_j\) se anula; la carga ponderada
por la orientación efectiva se anula asimismo al multiplicarse por
\(\epsilon_0\). El libro conserva la orientación inicial, la carga
acumulada y sus incrementos. Estas cantidades pertenecen
al cursor multiplicativo; los lectores de una trayectoria que
alterna ambas hojas actúan sobre el registro de esa trayectoria.

% END OWNER


La emisión finita, su acoplamiento de firmas y la reducción ternaria tienen
tipos sucesivos:
\[
 \mathcal S^2\xrightarrow{s\times e}
 \operatorname{im}s\times\operatorname{im}e
 \xrightarrow{G}\mathcal U_6
 \xrightarrow{\Psi_6}W_6\hookrightarrow\mathbb F_3^6.
\]
La sección~\ref{ii:sec:extension} demuestra la factorización de las
\(104\,976\) semillas en \(18\cdot26=468\) emisiones y las
\(243\) palabras visibles del ambiente de \(729\) elementos. Las
aplicaciones \(G\) y \(\Psi_6\) son las de
\eqref{ii:eq:acoplamiento} y~\eqref{ii:eq:psi-catalogo}; el registro conserva
las preimágenes aritméticas junto a esa publicación finita.

\subsubsection{Compatibilidad de cocientes y cilindros de refinamiento}

La prolongación compara estados que conservan simultáneamente residuos,
cocientes y fronteras. En la carta lineal de seis coordenadas del registro,
la matriz
\[
 A_H=\begin{pmatrix}
 2&2&2&1&2&1\\2&1&2&2&1&1\\1&0&0&1&0&2\\
 0&0&1&1&1&0\\2&2&2&0&2&1\\2&1&2&1&0&0
 \end{pmatrix},\qquad\det A_H=1,
\]
da la aplicación de división
\begin{equation}
 \mathcal H:\mathbb Z^6\longrightarrow
 \{0,1,2\}^6\times\mathbb Z^6,\qquad
 x\longmapsto(u,x^+),\quad
 y=xA_H,\quad u=[y]_3,\quad x^+=(y-u)/3.
 \label{ii:ii:tpk-hensel-local}
\end{equation}
Las filas se interpretan como vectores. El par \((u,x^+)\) conserva
\(y=u+3x^+\) y permite recuperar
\(x=(u+3x^+)A_H^{-1}\). La orientación ternaria de \(u\) se obtiene
después mediante \(\rho_3^{\rm or}\), conservando el cociente anterior.
Así se explicita el contenido reversible de esta carta local.

La carta cilíndrica relaciona un prefijo ternario de longitud \(L\), de
índice entero \(N\), con \(K\) tríadas decimales, de índice \(D\).
Los correspondientes cilindros semiabiertos tienen extremos racionales
\[
 I_3(N,L)=\left[\frac N{3^L},\frac{N+1}{3^L}\right),\qquad
 I_{1000}(D,K)=\left[\frac D{1000^K},\frac{D+1}{1000^K}\right).
\]
Definamos los enteros de frontera
\[
 a=N1000^K-D3^L,\qquad b=(N+1)1000^K-D3^L.
\]

\begin{lemma}[Actualización entera de los cilindros compatibles]
\label{ii:ii:tpk-cilindros-enteros}
Los cilindros anteriores se intersectan si y sólo si
\(a<3^L\) y \(b>0\). Al añadir un bloque ternario
\(u\in\{0,\ldots,728\}\), se tiene \(N'=729N+u\) y
\(L'=L+6\). Para \(\Delta K\in\{0,1\}\), las fronteras se
actualizan por
\[
 \begin{array}{ll}
 \Delta K=0:&a'=729a+u1000^K,\\[2pt]
 \Delta K=1:&D'=1000D+\delta,\quad
 a'=729000a+u1000^{K+1}-\delta\,729\,3^L,
 \end{array}
\]
donde \(\delta\in\{0,\ldots,999\}\) es la tríada de la
prolongación compatible. En ambos casos,
\(b'=a'+1000^{K+\Delta K}\).
\end{lemma}
\begin{proof}
Dos intervalos semiabiertos se intersectan precisamente cuando cada extremo
izquierdo es menor que el extremo derecho del otro. Multiplicar esas dos
desigualdades por \(3^L1000^K\) da \(a<3^L\), \(b>0\).
La sustitución de \(N'\), \(L'\) y \(D'\) en
\(a'=N'1000^{K+\Delta K}-D'3^{L'}\) produce las dos recurrencias.
Restar las definiciones de \(b'\) y \(a'\) da la última identidad.
\end{proof}

Estas operaciones actúan sobre prefijos enteros y sus restricciones
compatibles. La lectura arquimediana aparece al evaluar el sistema de
cilindros resultante. El dominio de prolongación incorpora además la
orientación, el acarreo y la firma conjunta de frontera de
\eqref{ii:ii:tpk-fibra}; la comparación de intervalos es una componente de
esa compatibilidad conjunta.

\subsubsection{Relaciones de extensión y conexión nonádica conjunta}

Para una ruta observable \(r:q\to q'\), la prolongación tiene tipo
\[
 \operatorname{Ext}_r\subseteq\mathcal F_q\times\mathcal F_{q'},
 \qquad
 \operatorname{Ext}_{r_2}\circ\operatorname{Ext}_{r_1}
 \subseteq\operatorname{Ext}_{r_2r_1}.
\]
Las identidades satisfacen
\(\Delta_{\mathcal F_q}\subseteq
\operatorname{Ext}_{\operatorname{id}_q}\).
Conserva los truncamientos, las hojas, las orientaciones, los acarreos y
las incidencias de frontera de las dos historias. Las identidades y las
concatenaciones admisibles forman la categoría de estados y prolongaciones
del TPK. Esta estructura relacional admite fibras locales con más de una
continuación y permite imponer conjuntamente su persistencia.

Sea \(\mathcal H_m\) el conjunto de historias admisibles de profundidad
\(m\), con truncamientos \(p_{n,m}\). La supervivencia se expresa por
\[
 \operatorname{Surv}_m^{(N)}=p_{N,m}(\mathcal H_N),\qquad
 \operatorname{Surv}_m=\bigcap_{N\geq m}
 \operatorname{Surv}_m^{(N)},\qquad
 X_\chi=\varprojlim_m\operatorname{Surv}_m^{(\chi)}.
\]
La sección compatible conserva simultáneamente todos sus prefijos. Las
relaciones locales \(\mathcal R_k\) de la conexión se componen en
\begin{equation}
 \Gamma_{9,k}=\mathcal R_{k+8}\circ\cdots\circ\mathcal R_k,\qquad
 \Gamma_{9,k}:\widehat X_k\dashrightarrow\widehat X_{k+9}.
 \label{ii:ii:tpk-holonomia}
\end{equation}
El dominio indicado es el de las historias compatibles. Sobre una historia
individualizada, las nueve relaciones representan su transporte efectivo;
el sistema conserva todas las prolongaciones admisibles antes de esa
individualización.

La ley de retorno se escribe
\[
 g(\Gamma_{9,k}x)=g(x),\qquad
 q_{\rm mem}(\Gamma_{9,k}x)=q_{\rm mem}(x)+1.
\]
La fase vuelve y la historia se prolonga. Los prefijos
\(w_6,w_{12},w_{18},w_{24},w_{30}\), la frontera \(R_{36}\) y las
cartas \(G_9\) son estados, extensiones o presentaciones de esta
construcción, con tipos distintos de la holonomía \(\Gamma_9\).
Su realización regional se desarrolla en la sección~\ref{ii:sec:generacion};
la composición de vueltas y sus representaciones están demostradas en
\ref{ii:mono-ampl:def-vuelta}--\ref{ii:mono-ampl:representacion}.

\subsubsection{Incidencia de fase y origen aritmético de los canales}

El selector \(M_{\rm ph}\) induce las clases
\[
 H_+=\{1,4,7\},\quad H_-=\{2,5,8\},\quad
 H_0=\{3,6,9\},\quad H_{\rm act}=H_+\sqcup H_-.
\]
El origen de fase fija el representante terminal
\(\star_{\rm ph}=9\) y el conjunto
\(O_{\rm ph}=D_9\setminus\{\star_{\rm ph}\}\).
Escribimos \(P_2(H_{\rm act})\) para los pares no ordenados de fases
activas y definimos
\[
 \mathcal F_6^{\rm ph}=H_{\rm act}\times P_2(H_{\rm act}),\qquad
 \mathcal F_8^{\rm ph}=O_{\rm ph}\times P_2(H_{\rm act}),\qquad
 \Theta_{\rm ph}=D_9\times P_2(H_{\rm act}).
\]

\begin{theorem}[Incidencia de la semicorona de fase]
\label{ii:ii:tpk-semicorona}
En el origen de ventana declarado,
\[
 |\mathcal F_6^{\rm ph}|=90,\qquad
 |\mathcal F_8^{\rm ph}|=120,\qquad
 |\Theta_{\rm ph}|=135.
\]
La frontera orientada, el retorno apuntado y el cierre lateral,
\[
 \partial_{\rm ph}=\Theta_{\rm ph}\times\{-1,+1\},\quad
 E_{\rm ph}=\partial_{\rm ph}\sqcup\{\infty\},\quad
 L_{\rm ph}=H_0\times P_2(H_{\rm act}),\quad
 P_{\rm ph}=\partial_{\rm ph}\sqcup L_{\rm ph},
\]
satisfacen
\[
 |\partial_{\rm ph}|=270,\quad |E_{\rm ph}|=271,\quad
 |L_{\rm ph}|=45,\quad |P_{\rm ph}|=315.
\]
La traslación del origen transporta por biyecciones todos estos conjuntos
y conserva sus incidencias.
\end{theorem}
\begin{proof}
Hay seis fases activas y \(\binom62=15\) pares no ordenados.
Los productos por \(6,8,9\) dan \(90,120,135\). La orientación
duplica el último conjunto, el retorno apuntado añade un elemento y las
tres fases neutrales aportan \(3\cdot15=45\) incidencias laterales.
La acción de \(C_9\) transporta simultáneamente el selector, el
representante terminal y cada factor de los productos cartesianos.
\end{proof}

La diferencia de cardinales activos es \(120-90=30\). El defecto
normalizado posee por ello la expresión exacta
\begin{equation}
 \frac{30}{120}-\frac{30}{90}=-\frac1{12}.
 \label{ii:ii:tpk-defecto-incidencia}
\end{equation}
Aquí \(90\) y \(120\) son cardinales de incidencias producidas por la
regla de fase. La posterior realización hexádico--octádica transporta esta
incidencia a sus objetos combinatorios. Las coordenadas angulares de esos
objetos se establecen mediante sus aplicaciones de realización.

\subsubsection{Registro temporal, extracción firmada y coordenadas reversibles}

Sea \(E_m=\{9(m-1)+1,\ldots,9m\}\). La restricción de una historia
a doce ventanas conserva, en cada una, su subruta \(\tau_m\),
orientación \(\sigma_m\), acarreo de frontera \(\kappa_m\) y registro
local \(\Lambda_m\):
\[
 \mathcal R_{12}(x)
 =\bigl((E_m,\tau_m,\sigma_m,\kappa_m,\Lambda_m)\bigr)_{m=1}^{12}.
\]
El mapa \(\mathcal R_{12}\) tiene como dominio las historias; el grupo
\(C_{12}\) tiene como elementos las clases de sector, y
\(\Theta_{108}\) conserva su calendario. Esta distinción mantiene
visibles los datos que cada realización conserva.

La extracción firmada de la sección~\ref{ii:sec:k-registro} factoriza por
\begin{equation}
 \mathcal X_{\rm TPK}^{\rm term,mem}
 \xrightarrow{\mathcal R_{12}}\mathcal R_{12}^{\rm mem}
 \xrightarrow{\mathcal E_{108}^{90,120}}
 \mathbb Z^{12}\times\mathbb Z^{12}\times\mathbb Z
 \xrightarrow{\Pi_H}\mathbb Z^{12}.
 \label{ii:ii:tpk-extraccion-tipada}
\end{equation}
El segundo mapa publica los canales
\((b^{90},b^{120},Q_{\rm TPK})\); el tercero construye el vector
firmado \(U\). La regla \(\varepsilon_9=M_{\rm ph}\) pertenece a la
selección de fase, mientras \(\mathcal E_{108}^{90,120}\) actúa sobre
el registro completo de canales. Sus dominios y sus salidas los distinguen.

Las ecuaciones~\eqref{ii:eq:k-sumas-armonicas} y
\eqref{ii:eq:k-bloques-firmados} dan la acción aritmética de \(\Pi_H\).
Después, la transformación \(\mathcal H_{12}\), construida por las
tres órbitas de paso tres, relaciona reversiblemente \(U\) con el
registro \(K\). Las diferencias cíclicas y la carga total proporcionan
lecturas posteriores del mismo registro. Así se conserva la procedencia
común de las coordenadas regionales y de la coordenada firmada antes de la
clausura entera que determina \(\alpha\).

\subsubsection{Transferencia reducida y memoria conservada por la dinámica}

Una vez construida la factorización de las emisiones,
\(\mathcal U_6\simeq\mathcal A_+\times\mathcal A_\times\), con
\(|\mathcal A_+|=18\) y \(|\mathcal A_\times|=26\), se dispone de
dos proyecciones sobre las funciones del catálogo:
\[
 (E_+f)(s,e)=\frac1{18}\sum_{s'\in\mathcal A_+}f(s',e),\qquad
 (E_\times f)(s,e)=\frac1{26}\sum_{e'\in\mathcal A_\times}f(s,e').
\]
Los alfabetos de firmas \(\mathcal A_+,\mathcal A_\times\) se
distinguen de los conjuntos de semillas de los dos cursores.
Cada una promedia la coordenada activa y conserva la otra. Con
\(P_{+1}=E_+\), \(P_{-1}=E_\times\), \(P_0=I\), el núcleo
de transferencia sobre \(\mathcal U_6\times C_9\) es
\[
 \mathcal K_{\rm ph}((u,r),(v,r+1))=P_{M_{\rm ph}(r)}(u,v).
\]
Las dos proyecciones conmutan y son idempotentes. Por tanto, su composición
durante una vuelta de fase satisface
\begin{equation}
 \mathcal K_{\rm ph}^{9}\big|_r=E_+E_\times,\qquad
 (E_+E_\times)(u,v)=\frac1{468}.
 \label{ii:ii:tpk-renovacion}
\end{equation}
Este operador es una publicación de memoria reducida sobre emisiones ya
construidas. La holonomía \(\Gamma_9\) conserva además las
coordenadas de la historia sobre la que se efectúa esa publicación.

En una realización isométrica \(U_{\Gamma_9}\) del transporte y para
una inclusión isométrica \(J\) de la coordenada observada, sus dos
componentes son
\[
 T=J^*U_{\Gamma_9}J,\qquad
 \eta=(I-JJ^*)U_{\Gamma_9}J.
\]
Se obtiene directamente
\[
 U_{\Gamma_9}J=JT+\eta,\qquad J^*\eta=0,\qquad
 I-T^*T=\eta^*\eta.
\]
La componente \(\eta\) se calcula a partir de la misma actualización
realizada y de la proyección declarada. La telescopía
\eqref{ii:eq:telescopia} conserva tanto las contribuciones expresadas en el
complemento ortogonal como el término terminal. Esta separación es la
base operatoria de la distinción posterior entre información registrada,
entropía de una observación y realización dimensional de la acción.
